% !TEX program = xelatex
% #1067 修复前后对比。数值由 measure.tex 分别在两个版本上实测得出, 此处静态排版。
\documentclass[11pt]{article}
\usepackage[margin=1.1cm,paperwidth=15.4cm,paperheight=8cm]{geometry}
\usepackage{xeCJK,xeCJKfntef,booktabs}
\setCJKmainfont{FandolSong-Regular.otf}
\pagestyle{empty}
\parindent=0pt
\begin{document}
{\bfseries \#1067：花括号分组内的中西文间距，收缩量是否回到外层}

\smallskip
判据：把正文压窄 2\,pt 时的 \texttt{\string\badness}。
有限值表示收缩量够用；\texttt{1000000} 表示不够、会溢出。

\medskip
\begin{tabular}{@{}llcc@{}}
\toprule
正文写法 & 自然宽 & 修复前 & 修复后 \\
\midrule
\verb|虚室 {hello} 生白|（无装饰命令，oracle） & 73.61\,pt & 55 & 55 \\
\verb|\CJKunderline{虚室 hello 生白}| & 73.61\,pt & 55 & 55 \\
\addlinespace
\verb|\CJKunderline{虚室 {hello} 生白}| & 73.61\,pt & \textbf{1000000} & \textbf{55} \\
\verb|\CJKunderline{虚室 \textbf{hello} 生白}| & 77.14\,pt & \textbf{1000000} & \textbf{55} \\
\verb|\CJKunderline{虚室 {{hello}} 生白}| & 73.61\,pt & \textbf{1000000} & \textbf{55} \\
\bottomrule
\end{tabular}

\medskip
三种带编组的写法都从「压不进」变为与 oracle 一致，自然宽度不变。
\end{document}
